\begin{table}[htbp]
\centering
\caption{Validation of the functions used for the trial examples and of the hazard formulas of Figure 1.}
\label{tab:functions}
\begin{threeparttable}
\scriptsize
\begin{tabular}{p{7cm}rrrrrr}
\toprule
Check & Comparisons & Largest difference & Tolerance & PASS & EXPLAINED & FAIL \\
\midrule
hand-made data vs Python & 28 & 4.9e-15 & 1e-12 & 28 & 0 & 0 \\
formulas vs Python & 5 & 3.3e-16 & 1e-12 & 5 & 0 & 0 \\
published value & 1 & 0.0e+00 &   0 & 1 & 0 & 0 \\
log-rank vs survival::survdiff (20 trials) & 8 & 2.7e-14 & 1e-08 & 8 & 0 & 0 \\
response Z vs stats::prop.test (20 trials) & 1 & 8.9e-16 & 1e-10 & 1 & 0 & 0 \\
win statistics vs BuyseTest & 25 & 4.4e-16 & 1e-06 & 25 & 0 & 0 \\
2-in-1 data & 2 & 0.0e+00 &   0 & 2 & 0 & 0 \\
2-in-1 statistics vs survival and stats (5 trials) & 5 & 2.5e-14 & 1e-08 & 5 & 0 & 0 \\
Figure 1 formulas & 5 & 2.6e-07 & 3.4e-05 & 5 & 0 & 0 \\
\bottomrule
\end{tabular}
\begin{tablenotes}
\footnotesize
\item Expected values of the hand-made data set were computed independently with loop-based Python code. Simulated data sets were compared with \texttt{survival::survdiff()}, \texttt{stats::prop.test()} and the \texttt{BuyseTest} package.
\end{tablenotes}
\end{threeparttable}
\end{table}
